\documentclass[10pt,a4paper]{amsart}
\usepackage[margin=19mm]{geometry}
\usepackage{amsmath,amssymb,mathtools}
\usepackage[scr=rsfs,cal=euler]{mathalfa}
\usepackage{xfrac}
\usepackage[unicode,hidelinks,bookmarks=false]{hyperref}
\newcommand{\ie}{i.e.\ }
\newcommand{\cf}{cf.\ }
\newcommand{\resp}{respectively\ }
\newcommand{\Subsection}{\subsection}
\providecommand{\nobreakdash}{\nobreak\hbox{-}}
\setlength{\emergencystretch}{2em}
\multlinegap0pt
\usepackage{mathtools}
\usepackage{amssymb}
\usepackage{dsfont}
\usepackage{bm}
\usepackage[shortlabels]{enumitem}
\usepackage{xcolor}
\usepackage{tikz}
\newtheorem{proposition}{Proposition}
\newtheorem{theorem}{Theorem}
\newtheorem{lemma}{Lemma}
\usepackage{cleveref}
\crefname{proposition}{Proposition}{Propositions}
\crefname{theorem}{Theorem}{Theorems}
\crefname{lemma}{Lemma}{Lemmas}
\newcommand{\CIID}{\mathcal{C}^{2}_{\mathrm{IID}}(\comp)}
\newcommand{\IID}{\mathrm{IID}}
\newcommand{\KIID}{K^{\mathrm{IID}}}
\newcommand{\LIID}{L^{\mathrm{IID}}}
\newcommand{\Qgam}{\mathcal{Q}_{\comp}}
\newcommand{\Val}{\normalfont\text{Val}}
\newcommand{\ac}[1]{#1}
\newcommand{\comp}{\gamma}
\newcommand{\ind}[1]{\mathds{1}_{\{#1\}}}
\title{Kernel Methods for Refined Prophet Inequalities}
\author{Patrick Loiseau, Mathieu Molina, Vianney Perchet, Sebastian Perez-Salazar, Victor Verdugo}
\date{Source: arXiv:2608.08662v1 (CC BY 4.0)}
\begin{document}
\maketitle

\noindent\emph{Source and licence.} Kernel Methods for Refined Prophet Inequalities. Version arXiv:2608.08662v1 (CC BY 4.0), distributed by arXiv under the Creative Commons Attribution 4.0 International licence, http://creativecommons.org/licenses/by/4.0/, as recorded in the arXiv licence metadata for this exact version and shown on the abstract page https://arxiv.org/abs/2608.08662v1. Full licence notice: \texttt{licenses/arxiv-2608.08662-CC-BY-4.0.txt}.\par\smallskip

\noindent\textit{Editorial scope.} Counted unit: the article's lemma lem:iid\_kernel\_validity (source0.tex lines 943--945). The article's complete original proof of it is delivered with it (source0.tex lines 3496--4144, one \texttt{proof} environment of 649 lines, which the article places away from the statement). No mathematics is written, repaired or re-derived: the statement and the proof are the article's own lines, in the article's own notation, and no retained line is substituted or re-flowed.  Retained uncounted as the source-local context the proof needs: lines 555--559; lines 822--838.  Nothing was omitted from any retained span, so the package carries no omission record.  No retained line contains a citation, so the wrapper carries no bibliography and no bibliography entry.  No retained line contains a figure environment, an image-inclusion command or drawing code, so no image is delivered and the package declares no asset dependency.  Editorial formatting only, in addition: an AMS \texttt{amsart} preamble that restates the article's own declarations and macros used by the retained lines, namely the environments \texttt{proposition}, \texttt{theorem}, \texttt{lemma} on separate counters, together with the standard packages those lines need.  The retained environments print in this document as Proposition 1, Theorem 1, Lemma 1 in order of appearance, whereas the article prints the same items with the numbering of its own full document.  The complete native source of the article as distributed in the arXiv e-print for this exact version is bundled unchanged as \texttt{sources/207179/source0.tex}.
\begin{proposition}\label{prop:iid_limit_kernel_game}
For every $\comp>0$, $\CIID
    =
    \Val(\Qgam,K^{\IID})$.
\end{proposition}

\begin{theorem}\label{thm:generic_variational_reduction}
Let $q \in (0,1)$. Suppose that $K$ is variationally regular, and suppose that $Q=1$ is not a minimizer of $L^K$ over $\Qgam$.  Then $V_K(q,s)>0$ for every $s\in[q,1)$, and minimizing $L^K$ in $Q$ is equivalent to minimizing over the one-parameter family $(Q_{q,s})_{s \in [q,1)}$, subject to non-negativity constraint.  Equivalently,
\begin{align*}\label{eq:generic_reduced_objective}
    &\inf_{Q\in\Qgam}\int_q^1  Q(u)K_q(u)\,du\\
    &=
    \inf_{\substack{
        s\in[q,1)\\
        Q_{q,s}\ge0
    }}
        \int_q^1K_q(u)\,du
        +
        \frac{\sqrt{\comp}}{\sqrt{V_K(q,s)}} \left(
        M_K(q,s)\int_q^1 K_q(u)\,du
        +
        \int_s^1(K_q(s)-K_q(u))K_q(u)\,du\right).
\end{align*}
\end{theorem}

\begin{lemma}\label{lem:iid_kernel_validity}
The \ac{IID} kernel $K^{\IID}$ is topologically admissible and variationally regular, and for every $Q\in\Qgam$, the map $q\mapsto L^{\IID}(Q,q)$ is quasi-concave on $[0,1]$.
\end{lemma}

\begin{proof}
We proceed in four steps.

\medskip

\noindent\emph{Step 1: duality and localization of the threshold parameter.}
By \Cref{prop:iid_limit_kernel_game}, the IID bounded-variance value is
the value of the limit-kernel game,
\begin{equation}\label{eq:iid_primal_variational}
    \CIID
    =
    \inf_{Q\in\Qgam}
    \sup_{q\in[0,1]}
    \LIID(Q,q),
\end{equation}
where, for $q\in(0,1)$,
\begin{equation*}
    \LIID(Q,q)
    =
    \frac{1-q}{-\log q}
    \int_q^1\frac{Q(u)}{u}\,du.
\end{equation*}
By \Cref{lem:iid_kernel_validity}, the IID kernel is topologically
admissible, and $q\mapsto\LIID(Q,q)$ is quasi-concave for every
$Q\in\Qgam$. Hence the strong-duality theorem gives
\begin{equation}\label{eq:iid_dual_variational}
    \CIID
    =
    \sup_{q\in[0,1]}
    \inf_{Q\in\Qgam}
    \LIID(Q,q).
\end{equation}

We first restrict the relevant range of $q$. The classical IID
single-threshold guarantee gives
\begin{equation*}
    \CIID\geq 1-\frac1e.
\end{equation*}
On the other hand, for every $q>1/e$, testing the inner problem on
the constant quantile $Q=1$ gives
\begin{equation*}
    \inf_{Q\in\Qgam}\LIID(Q,q)
    \leq
    \LIID(1,q)
    =
    1-q
    <
    1-\frac1e.
\end{equation*}
Therefore no $q>1/e$ can maximize the dual problem. At the lower
endpoint,
\begin{equation*}
    \inf_{Q\in\Qgam}\LIID(Q,0)
    =
    \inf_{Q\in\Qgam}Q(0^+)
    =
    0,
\end{equation*}
so $q=0$ is not an outer maximizer either.

We next verify that the constant quantile is not a minimizer for
$q\in(0,1/e)$. Fix such a $q$. Since $-\log q>1$ and
\begin{equation*}
    \lim_{r\uparrow1}\frac{-\log r}{1-r}=1,
\end{equation*}
we may choose $r\in(q,1)$ sufficiently close to $1$ so that
\begin{equation}\label{eq:iid_step_improvement_condition}
    \frac{-\log r}{1-r}<-\log q.
\end{equation}
Consider the normalized step quantile $S_r(u)=\frac{\ind{u>r}}{1-r}$.
It satisfies $\int_0^1S_r=1$, and
\begin{align*}
    \LIID(S_r,q)
    =
    \frac{1-q}{-\log q}
    \frac1{1-r}
    \int_r^1\frac{du}{u}
    =
    \frac{1-q}{-\log q}
    \frac{-\log r}{1-r}
    <
    1-q
    =
    \LIID(1,q).
\end{align*}
For $\varepsilon\in(0,1)$, define
\begin{equation*}
    Q_\varepsilon
    =
    (1-\varepsilon)1+\varepsilon S_r.
\end{equation*}
This function is nonnegative, nondecreasing, and has mean one.
Moreover,
\begin{align*}
    \int_0^1Q_\varepsilon(u)^2\,du
    &=
    (1-\varepsilon)^2
    +
    2\varepsilon(1-\varepsilon)\int_0^1S_r(u)\,du
    +
    \varepsilon^2\int_0^1S_r(u)^2\,du\\
    &=
    1+\varepsilon^2
    \left(\frac1{1-r}-1\right)\\
    &=
    1+\varepsilon^2\frac{r}{1-r}.
\end{align*}
Thus $Q_\varepsilon\in\Qgam$ for all sufficiently small
$\varepsilon>0$. By linearity,
\begin{equation*}
    \LIID(Q_\varepsilon,q)
    =
    (1-\varepsilon)\LIID(1,q)
    +
    \varepsilon\LIID(S_r,q)
    <
    \LIID(1,q).
\end{equation*}
Hence $Q=1$ is not a minimizer for any $q\in(0,1/e)$.


At $q=1/e$, the constant quantile is a minimizer. For every
nondecreasing $Q\in\Qgam$,
\begin{align*}
\int_{1/e}^1\frac{Q(u)}u\,du-\int_0^1Q(u)\,du
&=
\int_{1/e}^1Q(u)\left(\frac1u-1\right)du
-
\int_0^{1/e}Q(u)\,du.
\end{align*}
By monotonicity,
\begin{align*}
\int_{1/e}^1Q(u)\left(\frac1u-1\right)du
\geq
Q(1/e)\int_{1/e}^1\left(\frac1u-1\right)du
=
\frac{Q(1/e)}e,
\end{align*}
while
\begin{equation*}
\int_0^{1/e}Q(u)\,du
\leq
\frac{Q(1/e)}e.
\end{equation*}
Therefore,
\begin{equation*}
\int_{1/e}^1\frac{Q(u)}u\,du
\geq
\int_0^1Q(u)\,du
=
1.
\end{equation*}
It follows that
\begin{equation*}
\LIID(Q,1/e)
\geq
1-\frac1e
=
\LIID(1,1/e),
\end{equation*}
so $Q=1$ is a minimizer at $q=1/e$.



\medskip

\noindent\emph{Step 2: the Euler--Lagrange candidate solutions.}
Fix $q\in(0,1/e)$. The active IID kernel is
\begin{equation*}
    \KIID_q(u)
    =
    \frac{1-q}{-\log q}\frac1u,
    \qquad u\in[q,1].
\end{equation*}
It is strictly positive, strictly decreasing, and variationally
regular on $[q,1]$. By
\Cref{thm:generic_variational_reduction}, an inner minimizer may
therefore be taken in the contact form
\begin{equation}\label{eq:iid_contact_family}
    Q_s(u)
    =
    a_s+b_s\left(\frac1s-\frac1u\right)_+,
    \qquad
    s\in[q,1),
\end{equation}
where the positive factor $(1-q)/(-\log q)$ in the kernel has been
absorbed into $b_s$. Equivalently,
\begin{equation*}
    Q_s(u)
    =
    \begin{cases}
        a_s, & 0\leq u\leq s,\\[0.3em]
        a_s+b_s\left(\dfrac1s-\dfrac1u\right),
        & s<u\leq1.
    \end{cases}
\end{equation*}
The second-moment constraint is active, and $b_s>0$.

Let
\begin{equation*}
    h_s(u)
    =
    \left(\frac1s-\frac1u\right)_+.
\end{equation*}
We compute its first two moments. First,
\begin{align}
    H_1(s)
    \coloneqq
    \int_0^1h_s(u)\,du
    =
    \int_s^1
    \left(\frac1s-\frac1u\right)\,du
    =
    \frac{1-s}{s}+\log s
    =
    \frac1s-1+\log s.
    \label{eq:iid_H1}
\end{align}
Similarly,
\begin{align}
    H_2(s)
    \coloneqq
    \int_0^1h_s(u)^2\,du
    =
    \int_s^1
    \left(
        \frac1{s^2}
        -\frac{2}{su}
        +\frac1{u^2}
    \right)\,du
    =
    \frac{1-s}{s^2}
    +
    \frac{2\log s}{s}
    +
    \left(\frac1s-1\right)
    =
    \frac1{s^2}-1+\frac{2\log s}{s}.
    \label{eq:iid_H2}
\end{align}
Define
\begin{equation*}
    D(s)=H_2(s)-H_1(s)^2.
\end{equation*}
Using \eqref{eq:iid_H1} and \eqref{eq:iid_H2}, we obtain
\begin{align}
    D(s)
    =
    \frac{
        2s-2s^2+2s^2\log s-s^2\log^2s
    }{s^2}\notag
    =
    \frac{
        2(1-s)+s(2-\log s)\log s
    }{s}.
    \label{eq:iid_D}
\end{align}
Since $h_s$ is nonconstant for every $s<1$, $D(s)$ is its variance
under Lebesgue measure and is therefore strictly positive.

Writing $Q_s=a_s+b_sh_s$, the mean constraint gives
\begin{equation*}
    1
    =
    \int_0^1Q_s(u)\,du
    =
    a_s+b_sH_1(s),
\end{equation*}
and hence
\begin{equation}\label{eq:iid_a_from_b}
    a_s=1-b_sH_1(s).
\end{equation}
For the second moment,
\begin{align*}
    \int_0^1Q_s(u)^2\,du
    &=
    a_s^2+2a_sb_sH_1(s)+b_s^2H_2(s)\\
    &=
    1+b_s^2\left(H_2(s)-H_1(s)^2\right)\\
    &=
    1+b_s^2D(s).
\end{align*}
Since the second-moment constraint is active,
\begin{equation*}
    1+b_s^2D(s)=1+\comp.
\end{equation*}
Taking the positive root gives
\begin{align}
    b_s
    =
    \sqrt{\frac{\comp}{D(s)}}
    =
    \sqrt{
        \frac{\comp s}{
            2(1-s)+s(2-\log s)\log s
        }
    }.
    \label{eq:iid_b}
\end{align}
Consequently,
\begin{equation}\label{eq:iid_a}
    a_s
    =
    1-
    \sqrt{
        \frac{\comp s}{
            2(1-s)+s(2-\log s)\log s
        }
    }
    \left(
        \frac1s-1+\log s
    \right).
\end{equation}

\medskip

\noindent\emph{Step 3: the nonnegativity constraint.}
Since $b_s>0$ and $h_s\geq0$, the minimum value of $Q_s$ is $a_s$.
Therefore $Q_s\geq0$ if and only if
\begin{equation*}
    a_s\geq0.
\end{equation*}
Using \eqref{eq:iid_a_from_b}, this is equivalent to
\begin{align*}
    b_sH_1(s)\leq1
    \iff
    \comp H_1(s)^2\leq D(s)
    \iff
    \comp
    \leq
    \frac{D(s)}{H_1(s)^2}.
\end{align*}
Define
\begin{equation}\label{eq:iid_feasibility_ratio}
    R(s)
    =
    \frac{D(s)}{H_1(s)^2}
    =
    \frac{
        2s-2s^2+2s^2\log s-s^2\log^2s
    }{
        \left(1-s+s\log s\right)^2
    }.
\end{equation}

We verify that $R$ is strictly increasing on $(0,1)$. We obtain that
\begin{equation*}
    H_1'(s)
    =
    -\frac{1-s}{s^2},
    \qquad
    D'(s)
    =
    -\frac{2H_1(s)}s.
\end{equation*}
Therefore,
\begin{align}
    R'(s)
    &=
    \frac{
        D'(s)H_1(s)-2D(s)H_1'(s)
    }{
        H_1(s)^3
    }\notag\\
    &=
    \frac{
        2\left((1-s)D(s)-sH_1(s)^2\right)
    }{
        s^2H_1(s)^3
    }\notag\\
    &=
    \frac{
        2\left((1-s)^2-s\log^2s\right)
    }{
        s^3H_1(s)^3
    }.
    \label{eq:iid_R_derivative}
\end{align}
The numerator is strictly positive. Indeed, setting
$t=-\log s>0$, we have
\begin{equation*}
    \frac{1-s}{\sqrt{s}}
    =
    e^{t/2}-e^{-t/2}
    =
    2\sinh(t/2)
    >
    t
    =
    -\log s.
\end{equation*}
Squaring gives
\begin{equation*}
    (1-s)^2>s\log^2s.
\end{equation*}
Hence $R'(s)>0$.

Moreover,
\begin{equation*}
    R(s)\sim2s
    \qquad\text{as }s\downarrow0,
\end{equation*}
while
\begin{equation*}
    R(s)\sim\frac{4}{3(1-s)}
    \qquad\text{as }s\uparrow1.
\end{equation*}
Thus $R$ increases continuously from $0$ to $+\infty$. For every
$\comp>0$, there is therefore a unique
$s_{\min}(\comp)\in(0,1)$ such that
\begin{equation*}
    R(s_{\min}(\comp))=\comp,
\end{equation*}
that is,
\begin{equation}\label{eq:iid_smin_detailed}
    \comp
    =
    \frac{
        2s_{\min}-2s_{\min}^2
        +2s_{\min}^2\log s_{\min}
        -s_{\min}^2\log^2s_{\min}
    }{
        \left(
            1-s_{\min}+s_{\min}\log s_{\min}
        \right)^2
    }.
\end{equation}
Since $R$ is increasing, the nonnegativity condition is equivalent
to
\begin{equation*}
    s\geq s_{\min}(\comp).
\end{equation*}
Combining this with the contact-point condition $s\geq q$, the
feasible interval is
\begin{equation*}
    s\in[q\vee s_{\min}(\comp),1).
\end{equation*}

\medskip

\noindent\emph{Step 4: computation of the reduced objective.}
Let
\begin{equation*}
    c(q)=\frac{1-q}{-\log q}.
\end{equation*}
For $q\leq s$, substituting \eqref{eq:iid_contact_family} into the
IID payoff gives
\begin{align}
    \LIID(Q_s,q)
    &=
    c(q)\int_q^1\frac{Q_s(u)}u\,du\notag\\
    &=
    c(q)a_s\int_q^1\frac{du}{u}
    +
    c(q)b_s
    \int_s^1
    \left(
        \frac1s-\frac1u
    \right)
    \frac{du}{u}.
    \label{eq:iid_objective_split}
\end{align}
The first term is
\begin{equation*}
    c(q)a_s\int_q^1\frac{du}{u}
    =
    a_s(1-q).
\end{equation*}
For the second term,
\begin{align}
    J(s)
    &\coloneqq
    \int_s^1
    \left(
        \frac1s-\frac1u
    \right)
    \frac{du}{u}\notag\\
    &=
    \frac1s\int_s^1\frac{du}{u}
    -
    \int_s^1\frac{du}{u^2}\notag\\
    &=
    -\frac{\log s}{s}
    -
    \left(
        \frac1s-1
    \right)\notag\\
    &=
    1-\frac1s-\frac{\log s}{s}.
    \label{eq:iid_J}
\end{align}
Consequently,
\begin{align*}
    \LIID(Q_s,q)
    &=
    a_s(1-q)
    +
    \frac{1-q}{-\log q}\,b_sJ(s)\\
    &=
    (1-q)
    \left[
        1-b_sH_1(s)
        +
        \frac{b_s}{-\log q}J(s)
    \right]\\
    &=
    (1-q)
    \left[
        1+
        b_s
        \left(
            -H_1(s)
            +
            \frac{J(s)}{-\log q}
        \right)
    \right].
\end{align*}
Using
\begin{equation*}
    -H_1(s)
    =
    1-\frac1s-\log s
\end{equation*}
and \eqref{eq:iid_J}, we obtain
\begin{align}
    \LIID(Q_s,q)
    &=
    (1-q)
    \left[
        1+
        b_s
        \left(
            1-\frac1s-\log s
            +
            \frac1{-\log q}
            \left(
                1-\frac1s-\frac{\log s}{s}
            \right)
        \right)
    \right].
    \label{eq:iid_reduced_before_b}
\end{align}
Finally, substituting \eqref{eq:iid_b} gives
\begin{align}
    \psi(q,s) \!
    = \!
    (1-q)\! \!
    \left[
        1+
        \sqrt{
            \frac{\comp s}{
                2(1-s)+s(2-\log s)\log s
            }
        }
        \bigg(\!
            1-\frac1s-\log s
            +
            \frac1{-\log q}
            \big(
                1-\frac1s-\frac{\log s}{s}
            \big)\!
        \bigg)
    \!\right] \!.
    \label{eq:iid_psi_derived}
\end{align}

For every $q\in(0,1/e)$, the generic variational reduction and the
calculations above therefore give
\begin{equation}\label{eq:iid_fixed_q_reduction}
    \inf_{Q\in\Qgam}\LIID(Q,q)
    =
    \min_{s\in[q\vee s_{\min}(\comp),1]}
    \psi(q,s),
\end{equation}
where the value at $s=1$ is defined by continuous extension as $\psi(q,1)=1-q$. 
Indeed, as $s\uparrow1$,
\begin{equation*}
    H_1(s)=O((1-s)^2),
    \qquad
    J(s)=O((1-s)^2),
    \qquad
    D(s)=\Theta((1-s)^3),
\end{equation*}
and therefore
\begin{equation*}
    b_s\bigl(H_1(s)+|J(s)|\bigr)
    =
    O(\sqrt{1-s})
    \longrightarrow0.
\end{equation*}
This boundary corresponds to the constant quantile.

At $q=1/e$, the inner value is $1-1/e$, as shown in Step~1. The
reduced expression gives the same value. Indeed, since
$-\log(1/e)=1$, the term multiplying $b_s$ becomes
\begin{align*}
    &1-\frac1s-\log s
    +
    1-\frac1s-\frac{\log s}{s}\\
    &\qquad
    =
    \left(1+\frac1s\right)(-\log s)
    -
    2\left(\frac1s-1\right)
    \geq0,
\end{align*}
where the last inequality is
\begin{equation*}
    (1+t)\log t\geq2(t-1),
    \qquad t=\frac1s\geq1.
\end{equation*}
Thus the minimum is attained at $s=1$ and equals $1-1/e$.

At $q=0$, we use the continuous extension
\begin{equation*}
    \psi(0,s)
    =
    a_s
    =
    1-
    \sqrt{
        \frac{\comp s}{
            2(1-s)+s(2-\log s)\log s
        }
    }
    \left(
        \frac1s-1+\log s
    \right).
\end{equation*}
Its minimum over $s\geq s_{\min}(\comp)$ is zero, attained at
$s=s_{\min}(\comp)$, consistently with
\begin{equation*}
    \inf_{Q\in\Qgam}\LIID(Q,0)=0.
\end{equation*}

Combining the fixed-$q$ reduction with
\eqref{eq:iid_dual_variational} and the localization
$q\leq1/e$ yields
\begin{equation*}
    \CIID
    =
    \max_{q\in[0,1/e]}
    \min_{s\in[q\vee s_{\min}(\comp),1]}
    \psi(q,s),
\end{equation*}
with $\psi$ and $s_{\min}(\comp)$ given by
\eqref{eq:iid_psi_derived} and
\eqref{eq:iid_smin_detailed}, respectively.
\end{proof}

\end{document}
